\section{Why Both Judgments Are Fitting}

If every soul receives its irrevocable sentence immediately after death, a serious question naturally follows: why does the Creed still await a General Judgment at the end of the world? St.\ Thomas Aquinas places this exact objection at the head of his inquiry in the \textit{Summa Theologiae}, even quoting the ancient Septuagint reading of Nahum 1:9: ``God shall not judge the same thing a second time.''\endnote{St.\ Thomas Aquinas, \textit{Summa Theologiae} III, q.\ 59, a.\ 5, obj.\ 1--2; cf.\ \textit{Supplementum}, q.\ 88, a.\ 1, obj.\ 1--3.} If Lazarus already rests in Abraham's bosom and the rich man already suffers in flame, what remains to be judged?

Catholic theology answers without weakening either judgment. God does not reverse the Particular Judgment or punish the same sin twice; He brings to bodily, social, and cosmic completion what death settled for the separated soul alone.

In the \textit{Supplement}, St.\ Thomas anchors both judgments in the twofold reality of every human person:

\begin{quote}
Each man is both an individual person and a part of the whole human race: wherefore a twofold judgment is due to him. One, the particular judgment, is that to which he will be subjected after death, when he will receive according as he hath done in the body, not indeed entirely but only in part, since he will receive not in the body but only in the soul. The other judgment will be passed on him as a part of the human race.\endnote{\textit{Supplementum}, q.\ 88, a.\ 1, ad 1.}
\end{quote}

This duality mirrors God's own twofold action. In time, divine governance lets the good and the wicked remain mingled so that the just may be tried and the wicked called to conversion; the Particular Judgment belongs to that governance of individuals. At the close of history, when no further reason remains for wheat and cockle to grow together, the General Judgment corresponds to the original creation of the cosmos, bringing the entire world to its ultimate complement ``by each one receiving finally his own personal due.''\endnote{\textit{Supplementum}, q.\ 88, a.\ 1, corpus.}

In the \textit{Tertia Pars}, St.\ Thomas adds a sharp metaphysical principle: ``Judgment cannot be passed perfectly upon any changeable subject before its consummation.'' Although death fixes the soul's own essential state, ``man's temporal life in itself ends with death, still it continues dependent in a measure on what comes after it in the future.''\endnote{Aquinas, \textit{ST} III, q.\ 59, a.\ 5, corpus.} St.\ Thomas lists five threads that outlive a person's final breath: human memory and reputation, often false to the truth; one's children; the enduring historical effects of one's works and teaching---``just as from the deceit of Arius and other false leaders unbelief continues to flourish down to the close of the world; and even until then faith will continue to derive its progress from the preaching of the apostles''; the mortal body turning to dust until the resurrection; and the temporal things on which a heart was set. Therefore, ``there must be a final Judgment at the last day, in which everything concerning every man in every respect shall be perfectly and publicly judged.''\endnote{Aquinas, \textit{ST} III, q.\ 59, a.\ 5, corpus and ad 1.}

The \textit{Roman Catechism} distills these truths into four classical reasons that pastors are enjoined to teach the faithful.\endnote{\textit{Catechism of the Council of Trent}, Part I, Creed Article VII, trans.\ Donovan, pp.\ 62--63.}

First stands the enduring ripple of good and evil example. Parents, teachers, writers, rulers, and neighbors leave behind chains of influence that end only when history ends; because ``the good or bad influence of example, affecting as it does, the conduct of many, is to terminate only with this world; justice demands that, in order to form a proper estimate of the good or bad actions of all, a general judgment should take place.'' While separated souls acquire no new essential merit or demerit after death, the public unveiling of every soul blessed by their fidelity or harmed by their scandal belongs to the complete accounting of their lives.\endnote{Cf.\ Aquinas, \textit{ST} III, q.\ 59, a.\ 5, ad 1.}

Second stands the public vindication of slandered virtue and unmasking of hypocrisy. History frequently calumniates the innocent while praising the corrupt. The Book of Wisdom foretells the great reversal:

\begin{quote}
Then shall the just stand with great constancy against those that have afflicted them, and taken away their labours. These seeing it, shall be troubled with terrible fear \dots\ saying within themselves, repenting, and groaning for anguish of spirit: These are they, whom we had sometime in derision, and for a parable of reproach. We fools esteemed their life madness, and their end without honour. Behold, how they are numbered among the children of God, and their lot is among the saints. (Wis 5:1--5)
\end{quote}

Christ promises that ``nothing is covered that shall not be revealed: nor hid, that shall not be known'' (Matt 10:26), and St.\ Paul bids Christians await the Lord, ``who both will bring to light the hidden things of darkness and will make manifest the counsels of the hearts. And then shall every man have praise from God'' (1 Cor 4:5; cf.\ Eccles 12:14). Thus, says the \textit{Roman Catechism}, divine justice restores to the virtuous ``in the presence and with the suffrage of a congregated world, the good name of which they had been unjustly deprived before men.''\endnote{\textit{Catechism of the Council of Trent}, Part I, Creed Article VII, trans.\ Donovan, p.\ 62.}

Third stands the unity of body and soul. Man is not a disembodied spirit; the body cooperated as the soul's instrument in every work of charity and every sin. At the Particular Judgment, only the separated soul receives its due. Only through the general resurrection and General Judgment can ``the body, therefore, \dots\ participate with the soul in the eternal rewards of virtue, or the everlasting punishments of vice,'' fulfilling St.\ Paul's word that each must receive ``the proper things of the body''---\textit{propria corporis}---``according as he hath done, whether it be good or evil'' (2 Cor 5:10).\endnote{\textit{Catechism of the Council of Trent}, Part I, Creed Article VII, trans.\ Donovan, p.\ 62; Aquinas, \textit{ST} III, q.\ 59, a.\ 5, ad 3; \textit{Supplementum}, q.\ 88, a.\ 1, ad 1--2.}

Fourth stands the public vindication of Divine Providence. Even the Psalmist nearly stumbled at seeing sinners prosper while the innocent suffered (Ps 72:2--3, 12--14). In Book XX of \textit{The City of God}, St.\ Augustine surveys this mingled web of human affairs---the poor righteous man and the feasting reprobate, the innocent condemned by false evidence and the ruffian crowned with honors---noting that because earthly good and ill fall indifferently on good and bad alike, God trains us to despise transitory things and seek what belongs to the good alone. Yet hidden justice will not stay hidden forever:

\begin{quote}
But when we shall have come to that judgment, the date of which is called peculiarly the day of judgment, and sometimes the day of the Lord, we shall then recognise the justice of all God's judgments, not only of such as shall then be pronounced, but of all which take effect from the beginning, or may take effect before that time. And in that day we shall also recognise with what justice so many, or almost all, the just judgments of God in the present life defy the scrutiny of human sense or insight, though in this matter it is not concealed from pious minds that what is concealed is just.\endnote{St.\ Augustine, \textit{The City of God} XX.2, trans.\ Dods (1871).}
\end{quote}

Then every murmur against Providence falls silent in ``the day of wrath and revelation of the just judgment of God \dots\ For there is no respect of persons with God'' (Rom 2:5, 11).\endnote{\textit{Catechism of the Council of Trent}, Part I, Creed Article VII, trans.\ Donovan, pp.\ 62--63.}

Crowning all four reasons stands the public glory of Christ Himself. As the \textit{Baltimore Catechism No.\ 3} teaches, a General Judgment is needed ``especially that Christ Our Lord may receive from the whole world the honor denied Him at His first coming, and that all may be forced to acknowledge Him their God and Redeemer.''\endnote{\textit{Baltimore Catechism No.\ 3} (1921), Q.\ 1389, p.\ 309.} As the \textit{Catechism of the Catholic Church} teaches, that Day will disclose the ultimate meaning of creation and of the whole economy of salvation, laying bare how divine Providence guided all things toward their end and showing that God's justice overcomes every human wrong and His love proves stronger than death.\endnote{\textit{Catechism of the Catholic Church}, \S 1040.} The two judgments stand together like the private sealing of a verdict and its solemn promulgation before the universe.
